\documentclass[10pt]{article}

\usepackage[utf8]{inputenc}

\usepackage[T1]{fontenc}

\usepackage{amsmath}

\usepackage{amsfonts}

\usepackage{amssymb}

\usepackage[version=4]{mhchem}

\usepackage{stmaryrd}



\begin{document}

Здесь $L-$ характерный размер, $v-$ характерная скорость, $v=(\eta / \rho)-$ кинематич. вязкость. Диффузией магнитного поля можно пренебречь, если велико магнитное число Рейнольдса:



\[

\operatorname{Re}_{m}=\frac{L v}{v_{m}} .

\]



Наконец, теплопроводностью можно пренебречь, если велико число Пекле:



\[

\mathrm{Pe}=\frac{L v}{\chi}

\]



где $\chi=\frac{\chi}{\rho c_{p}} \quad$ коэффициент температуропроводности, $c_{p}-$ удельная теплоемкость при постоянном давлении.



Равновесные конфигурации. Такие конфигурации (м а г н и т ны е ло в уш к и) представляют интерес для осуществления управляемых термоядерных реакций. При этом магнитные силовые линии целиком лежат на поверхностях



\[

p(x, y, z)=\text { const, }

\]



называемых поэтому магнитными поверхностям и. Простейшая равновесная конфигурация аксиально симметрична. В этом случае величины в М.г.у. зависят только от двух цилиндрич. координат $r$ и $z$ и не зависят от азимута $\vartheta$. Аксиально симметричное магнитное поле называется то ро и даль ны м, если отлична от нуля только его азимутальная компонента $B_{\vartheta}$. Если же отличны от нуля только компоненты $B_{r}$ и $B_{z}$, то поле называется полои да л ьны м. Аксиально-симметричные равновесные конфигурации описываются уравне н и м Грэда-Шафра но в а



\[

\frac{\partial^{2} \psi}{\partial r^{2}}-\frac{1}{r} \frac{\partial \psi}{\partial r}+\frac{\partial^{2} \psi}{\partial z^{2}}+16 \pi^{3} r^{2} \frac{d p}{d \psi}+\frac{8 \pi^{2}}{c^{2}} \frac{d\left(I^{2}\right)}{d \psi}=0,

\]



где $\psi$ - магнитный поток через круг радиусом $r$, находящийся на высоте $z$ :



\[

\psi(r, z)=\int_{0}^{r} B(z, \rho) 2 \pi \rho d \rho,

\]



и $I$ - ток, протекающий через тот же круг



\[

I(r, z)=\int_{0}^{r} j_{z}(\rho, z) 2 \pi \rho d \rho

\]



$c$ - скорость света.



Функции $\psi(r, z)$ и $I(r, z)$ определяют магнитное поле:



\[

B_{z}=\frac{1}{2 \pi r} \frac{\partial \psi}{\partial r} ; \quad B_{r}=\frac{1}{2 \pi r} \frac{\partial \psi}{\partial z} ; \quad B_{\vartheta}=\frac{2 I}{c r}

\]



и плотность тока $j$ согласно уравнению (7). Величины $\psi$ и $I$ постоянны на магнитных поверхностях. Задавая $I$ как функцию от $\psi$, можно получить много простых аналитич. решений уравнения (9), что позволяет исследовать широкий класс равновесных плазменных конфигураций. Особую роль для управляемых термоядерных реакций играют за м кнутые магнитные ловушки, то есть конфигурации, топологически эквивалентные тору. Если плазма удерживается внешним магнитным полем, то тороидальную ловушку называют с те л л а р а т о ро м. Если же удержание плазмы осуществляется магнитным полем азимутального тока, текущего внутри плазмы, то ловушка называется то ка ма ком.



Стационарные течения. Простейшим магнитогидродинамич. течением в канале, служащим моделью многих инженерных приложений, является течение вязкой несжимаемой жидкости с учетом конечной электропроводности между двумя параллельными плоскостями $x= \pm l$ (течен и Гартмана). При этом



\[

v_{z}(x)=-\frac{l^{2}}{\eta H a \operatorname{sh} H a} \frac{\partial p}{\partial z}\left(\operatorname{ch} H a-\operatorname{ch} \frac{H a \cdot x}{l}\right), \quad v_{x}=v_{y}=0,

\]



где $H a=l V_{A} / \sqrt{v v_{m}}, \frac{\partial p}{\partial z}-$ градиент давления.



Турбулентность. В случае неустойчивости равновесного состояния в плазме возбуждается большое число волн, имеющих случайный характер. Такое состояние плазмы назы- вается турбулентным. Напротив, упорядоченное движение спокойной плазмы, при к-ром она движется как бы отдельными слоями, обладающими различными, но определенными скоростями, называется л а м и на р ным.



Слабая турбулентность. При малом уровне турбулентности, когда нелинейные члены в М. г. у. уже следует учитывать, но они все еще невелики, турбулентное состояние можно рассматривать как газ слабо взаимодействующих волн - плазмонов, имеющих случайные фазы. Такую турбулентность называют слабой. В случае $\beta \equiv \frac{8 \pi p}{B^{2}} \ll 1$ наиболее сушественны следующие процессы взаимодействия волн:



\[

A \rightleftarrows A+S ; \quad A \rightleftarrows F+S ; \quad F \rightleftarrows F+S ; \quad F \rightleftarrows A+S .

\]



Здесь буквами $A, F, S$ обозначены соответственно альфвеновская, быстрая и медленная магнитозвуковые волны.



Аномальные процессы переноса. В турбулентной среде коэффициенты переноса становятся аномально большими. В частности, аномальная электропроводность



\[

\sigma^{(T)}=3 c^{2} /\left(4 \pi \tau_{c}\left\langle\tilde{y}^{2}\right\rangle\right),

\]



где $c-$ скорость света, $\left\langle\widetilde{v}^{2}\right\rangle-$ средний квадрат флуктуации скорости, $\tau_{c}-$ время корреляции флуктуаший.



Гидромагнитное динамо. Под г и д ро м а гн и т ны м д и н а м о понимают механизм, способный генерировать магнитное поле в движущейся проводящей среде. Основой гидромагнитного динамо является растяжение силовых линий «затравочного» магнитного поля. При вращении звезды различные ее точки движутся с различной скоростью. Благодаря закону вмороженности это приводит к растяжению магнитных силовых линий. Что же касается затравочного поля, то оно образуется током, возникающим вследствие различия вязкости электронов и ионов. Однако теория ламинарного динамо не объясняет возникновения полоидального магнитного поля. Последнее связано с турбулентностью. При турбулентном движении рост магнитного поля пропорционален его ротору ( $\alpha$-эффект)



Здесь



\[

\frac{\partial\langle\boldsymbol{B}\rangle}{\partial t}=\alpha \operatorname{rot}\langle\boldsymbol{B}\rangle

\]



\[

\alpha=-\frac{1}{3}\langle\tilde{v}(t) \operatorname{rot} \tilde{v}(t)\rangle \tau_{c}

\]



Jum.: [1] Альфвен Г., Фельтхам мар К. Г., Космическая электродинамика, пер. с англ., 2 изд., М., 1967; [2] Ку ли ко в -

с к и й А. Г., Л ю б и м о в Г.А., Магнитная гидродинамика, М., 1962; [3] Л а нда у Л. Д., Л и фш и ц Е. М., Электродинамика сплошных сред, М., 1982; [4] Поло в и н Р. В., Дем мшки й В. П., Основы магнитной гидродинамики, М., 1987. В.П. Демуцкий, Р. В. Половин. МАТНМТНЫЙ ТОК - см. Ток. МАГНИТНЫЙ ФОРМФАКТОР - см. Формфактор. МАГНИТОГИДРОДИНАМИЧЕСКАЯ ТУРБУЛЕНТНОСТЬ - см. Гидромагнитная турбулентность.



МАГНИТОГИДРОДИНАМИЧЕСКИЕ ВОЛНЫ - вОЛНЫ, описываемые магнитной гидродинамики уравнениями. Существуют четыре типа волн.



\begin{enumerate}

  \item Альфвеновская волна, движущаяся со скоростью

\end{enumerate}



\[

V_{A n}=\frac{B_{n}}{\sqrt{4 \pi \rho}}

\]



\begin{enumerate}

  \setcounter{enumi}{1}

  \item Два типа магнит о в уко в ы в о л н, движущихся со скоростями

\end{enumerate}



\[

V=\sqrt{\frac{1}{2}\left[V_{A}^{2}+a^{2} \pm \sqrt{\left(V_{A}^{2}+a^{2}\right)^{2}-4 V_{A n}^{2} a^{2}}\right]} .

\]



Здесь $a-$ скорость звука, $V_{A}=|\boldsymbol{B}| / \sqrt{4 \pi \rho}$, знак + соответствует скорости быстрой магнитозвуковой волны $V_{F}$, знак - соответствует скорости медленной магнитозвуковой волны $V_{S}$.





\end{document}